\documentclass[aspectratio=169]{beamer}
\usetheme{default}
\usecolortheme{default}
\setbeamertemplate{navigation symbols}{}
\setbeamertemplate{footline}{%
    \hfill\textcolor{gray}{\footnotesize\insertframenumber}\hspace{1em}\vspace{0.5em}
}

\definecolor{fillblue}{RGB}{14, 52, 75}
\definecolor{fillaccent}{RGB}{183, 1, 0}
\definecolor{filllight}{RGB}{250, 240, 240}

\setbeamercolor{frametitle}{fg=fillblue}
\setbeamercolor{title}{fg=white}
\setbeamercolor{subtitle}{fg=white!80}
\setbeamercolor{author}{fg=white!90}
\setbeamercolor{itemize item}{fg=fillaccent}
\setbeamercolor{itemize subitem}{fg=fillblue}

\setbeamertemplate{frametitle}{%
    \vspace{0.6em}%
    {\large\bfseries\color{fillblue}\insertframetitle}%
    \par\vspace{0.15em}%
    {\color{fillaccent}\rule{0.12\textwidth}{1.5pt}}%
    \par\vspace{0.3em}%
}

\setbeamertemplate{title page}{%
    \vfill
    \begin{center}
        {\usebeamerfont{title}\usebeamercolor[fg]{title}\inserttitle\par}%
        \vspace{0.3em}%
        {\color{fillaccent}\rule{0.6\textwidth}{0.8pt}\par}%
        \vspace{1.4em}%
        {\usebeamerfont{subtitle}\usebeamercolor[fg]{subtitle}\insertsubtitle\par}%
        \vspace{1.0em}%
        {\usebeamerfont{author}\usebeamercolor[fg]{author}\insertauthor\par}%
    \end{center}
    \vfill
}

\usepackage[T1]{fontenc}
\usepackage[utf8]{inputenc}
\usepackage[naustrian]{babel}
\usepackage{lmodern}
\usepackage{graphicx}
\graphicspath{{../shared/images/}}
\usepackage{tikz}
\usetikzlibrary{positioning, arrows.meta}
\usepackage{booktabs}
\usepackage{colortbl}
\usepackage{array}

\newcommand{\plan}{\cellcolor{fillaccent!35}}
\newcommand{\darkbackground}{%
    \usebackgroundtemplate{%
        \begin{tikzpicture}[remember picture, overlay]
            \fill[fillblue] (current page.south west) rectangle (current page.north east);
            \fill[fillaccent] (current page.south west) rectangle ([yshift=0.5cm]current page.south east);
        \end{tikzpicture}%
    }%
}


% -----------------------------------------------
% Metadaten – hier anpassen
% -----------------------------------------------
\title{Titel der Präsentation}
\subtitle{Untertitel}
\author{Vorname Nachname}
\institute{}
\date{}

% -----------------------------------------------
\begin{document}
% -----------------------------------------------

% ============================
% FOLIE 1: Deckblatt
% ============================
{
\setbeamertemplate{footline}{}
\darkbackground
\begin{frame}
    \titlepage
\end{frame}
}

% ============================
% FOLIE 2: Aufzählung
% ============================
\begin{frame}{Folientitel}
    \begin{itemize}
        \setlength{\itemsep}{0.8em}
        \item Erster Punkt
        \item Zweiter Punkt mit \textbf{Hervorhebung}
        \item Dritter Punkt
    \end{itemize}
\end{frame}

% ============================
% FOLIE 3: Zwei Spalten
% ============================
\begin{frame}{Zwei Spalten}
    \begin{columns}[T]
        \begin{column}{0.48\textwidth}
            \textbf{\color{fillblue} Linke Spalte}\\[0.5em]
            \begin{itemize}
                \item Punkt A
                \item Punkt B
                \item Punkt C
            \end{itemize}
        \end{column}
        \begin{column}{0.48\textwidth}
            \textbf{\color{fillblue} Rechte Spalte}\\[0.5em]
            \begin{itemize}
                \item Punkt D
                \item Punkt E
                \item Punkt F
            \end{itemize}
        \end{column}
    \end{columns}
\end{frame}

% ============================
% FOLIE 4: Hervorgehobene Frage / Aussage
% ============================
\begin{frame}{Zentrale Frage}
    \vspace{0.6em}
    \begin{center}
        \setlength{\fboxrule}{1.5pt}%
        \fcolorbox{fillaccent}{white}{\parbox{0.88\textwidth}{%
            \vspace{0.4em}
            \centering
            {\large\color{fillblue}
            Hier steht eine zentrale Frage oder\\
            eine wichtige Aussage die hervorgehoben werden soll.}
            \vspace{0.4em}
        }}
    \end{center}
    \vspace{0.7em}
    \textbf{Unterpunkte}\\[0.4em]
    \begin{itemize}
        \setlength{\itemsep}{0.3em}
        \item Erster Unterpunkt
        \item Zweiter Unterpunkt
        \item Dritter Unterpunkt
    \end{itemize}
\end{frame}

% ============================
% FOLIE 5: Zeitplan-Tabelle
% ============================
\begin{frame}{Zeitplan}
    \vspace{0.8em}
    \large
    \setlength{\tabcolsep}{18pt}
    \renewcommand{\arraystretch}{1.8}
    \begin{center}
        \begin{tabular}{lccc}
            \toprule
            \textbf{Arbeitspaket} & \textbf{Monat 1} & \textbf{Monat 2} & \textbf{Monat 3} \\
            \midrule
            Arbeitspaket 1  & \plan &       &       \\
            Arbeitspaket 2  & \plan & \plan &       \\
            Arbeitspaket 3  &       & \plan &       \\
            Arbeitspaket 4  &       &       & \plan \\
            \bottomrule
        \end{tabular}
    \end{center}
\end{frame}

% ============================
% FOLIE 6: Abschluss
% ============================
{
\setbeamertemplate{footline}{}
\darkbackground
\begin{frame}
    \vspace{4em}
    \begin{center}
        {\LARGE\bfseries\color{white} Fragen?}
    \end{center}
\end{frame}
}

\end{document}